\section{Novel Predictions: Quantitative Tests of the Framework}
\label{sec:predictions}

%%%%%%%%%%%%%%%%%%%%%%%%%%%%%%%%%%%%%%%%%%%%%%%%%%%%%%%%%%%%%%%%%%%%%%%%%%%%%%%
% EXTENDED VERSION: Each prediction as comprehensive subsection
% Integrating Γ-matrix coefficients and LLM-MC methodology from Chapter 10
%%%%%%%%%%%%%%%%%%%%%%%%%%%%%%%%%%%%%%%%%%%%%%%%%%%%%%%%%%%%%%%%%%%%%%%%%%%%%%%

A unified framework must generate predictions that no constituent theory 
produces alone. This chapter develops eight novel predictions that emerge 
from the interaction of complementarity, context, and coherence. Each 
prediction is:
\begin{enumerate}
  \item Derived from the formal $(C, \Psi, K, Q)$ structure
  \item Quantified using the $\Gamma$-matrix estimated in Chapter~10
  \item Empirically testable with specified validation criteria
  \item Policy-relevant with concrete implications
\end{enumerate}

The predictions span trade policy, financial contagion, technological disruption,
identity economics, development traps, climate governance, and mechanism design---
demonstrating the framework's scope while maintaining analytical precision.

%%%%%%%%%%%%%%%%%%%%%%%%%%%%%%%%%%%%%%%%%%%%%%%%%%%%%%%%%%%%%%%%%%%%%%%%%%%%%%%
\subsection{What Makes a Prediction ``Novel''?}
%%%%%%%%%%%%%%%%%%%%%%%%%%%%%%%%%%%%%%%%%%%%%%%%%%%%%%%%%%%%%%%%%%%%%%%%%%%%%%%

\subsubsection{Criteria for Novelty}

A prediction qualifies as novel if it satisfies four conditions:

\begin{enumerate}
  \item \textbf{Emergence:} No single constituent theory generates it alone
  \item \textbf{Interaction:} It emerges from the \emph{combination} of framework elements
  \item \textbf{Testability:} It specifies observable implications that could be falsified
  \item \textbf{Surprise:} It differs from standard intuitions or existing predictions
\end{enumerate}

\subsubsection{The Framework's Unique Combinatorial Structure}

The $(C, \Psi, K, Q)$ framework uniquely combines elements that existing theories treat separately:

\begin{center}
\renewcommand{\arraystretch}{1.3}
\begin{tabular}{lll}
\hline
\textbf{Element} & \textbf{Existing Treatment} & \textbf{Framework Integration} \\
\hline
Complementarities & Game theory (strategic) & Cross-domain via $C$-matrix \\
Context & ``Control variables'' & Endogenous $\Psi$-dynamics \\
Coherence & Equilibrium (yes/no) & Continuous $K \in [0,1]$ \\
Welfare & GDP or utility & Six FEPSDE dimensions \\
Heterogeneity & Fixed effects & $\Gamma$-matrix interactions \\
\hline
\end{tabular}
\end{center}

\subsubsection{The Quantitative Foundation: The $\Gamma$-Matrix}

Chapter~10 established the $\Gamma$-matrix linking eight context dimensions ($\Psi$) 
to six welfare dimensions (FEPSDE). Key coefficients driving predictions include:

\begin{center}
\small
\begin{tabular}{lcl}
\hline
\textbf{Coefficient} & \textbf{Value (SE)} & \textbf{Interpretation} \\
\hline
$\gamma_{S,M} = -0.18$ & (0.04) & Market integration reduces social welfare \\
$\gamma_{Eco,M} = -0.21$ & (0.04) & Market integration reduces ecological welfare \\
$\gamma_{E,S} = +0.47$ & (0.06) & Social trust strongly predicts emotional wellbeing \\
$\gamma_{F,I} = +0.31$ & (0.05) & Institutions enable financial welfare \\
$\gamma_{Eco,E} = -0.19$ & (0.04) & Economic development reduces ecological welfare \\
$\gamma_{D,K} = +0.39$ & (0.05) & Information access enables digital welfare \\
\hline
\end{tabular}
\end{center}

These coefficients, estimated via LLM Monte Carlo (Appendix~AN) and validated 
against empirical benchmarks (MAE = 0.034), enable quantitative predictions 
unavailable to purely qualitative frameworks.

%%%%%%%%%%%%%%%%%%%%%%%%%%%%%%%%%%%%%%%%%%%%%%%%%%%%%%%%%%%%%%%%%%%%%%%%%%%%%%%
\subsection{Prediction 1: Trade War Dynamics and Hysteresis}
\label{subsec:pred-tradewar}
%%%%%%%%%%%%%%%%%%%%%%%%%%%%%%%%%%%%%%%%%%%%%%%%%%%%%%%%%%%%%%%%%%%%%%%%%%%%%%%

\subsubsection{The Policy Context}

Trade conflicts have intensified since 2018, with tariffs imposed on hundreds 
of billions of dollars of goods between major economies. The 2025 tariff 
escalations continue this pattern. Standard trade theory predicts welfare 
losses from tariffs---but struggles to explain their political popularity 
and persistent effects.

\subsubsection{Standard Economic Prediction}

Classical trade theory \citep{krugman1980} predicts:
\begin{itemize}
  \item Tariffs create deadweight loss: both countries lose
  \item Effects are primarily financial ($U^F$)
  \item Removal of tariffs restores pre-war equilibrium
  \item Political support for tariffs is ``irrational'' or explained by special interests
\end{itemize}

The prediction is symmetric and reversible: tariffs on, welfare down; tariffs off, welfare restored.

\subsubsection{Framework Prediction: Multi-Dimensional Effects}

The $(C, \Psi, Q)$ framework predicts a fundamentally different pattern 
by incorporating the $\Gamma$-matrix structure.

\paragraph{Phase 1: Immediate Effects ($t = 0$--$1$).}

Tariffs change market integration ($\Psi_M$). Using the $\Gamma$-matrix:

\begin{equation}
\Delta U_d = \gamma_{d,M} \cdot \Delta \Psi_M \quad \text{for each FEPSDE dimension } d
\end{equation}

\begin{center}
\renewcommand{\arraystretch}{1.3}
\begin{tabular}{lccc}
\hline
\textbf{Dimension} & $\gamma_{d,M}$ & $\Delta \Psi_M$ (tariffs) & $\Delta U_d$ \\
\hline
Financial ($U^F$) & $+0.28$ & $-0.3$ & $-0.084$ \\
Emotional ($U^E$) & $+0.02$ & $-0.3$ & $-0.006$ \\
Physical ($U^P$) & $+0.03$ & $-0.3$ & $-0.009$ \\
Social ($U^S$) & $-0.18$ & $-0.3$ & $\mathbf{+0.054}$ \\
Digital ($U^D$) & $+0.24$ & $-0.3$ & $-0.072$ \\
Ecological ($U^{Eco}$) & $-0.21$ & $-0.3$ & $\mathbf{+0.063}$ \\
\hline
\end{tabular}
\end{center}

\textbf{Key insight:} While $U^F$ and $U^D$ decline, $U^S$ and $U^{Eco}$ \emph{increase}.
The negative $\gamma_{S,M}$ coefficient means reduced market integration 
\emph{improves} social welfare---capturing the community cohesion effects 
documented by \citet{autor2013china}.

\paragraph{Phase 2: Identity and Coherence Effects ($t = 1$--$2$).}

Tariffs are not merely economic policy; they signal identity positions.
Let $U^{IDN}$ denote identity utility (Chapter~8):

\begin{equation}
U^{IDN} = \omega_{IDN} \cdot \mathbf{1}[\text{policy aligns with group identity}]
\end{equation}

For voters with ``producer'' or ``nationalist'' identity:
\begin{itemize}
  \item Tariffs signal ``we fight for our workers'' $\Rightarrow$ $U^{IDN} \uparrow$
  \item Financial losses framed as ``sacrifice for the nation'' $\Rightarrow$ $\lambda_F \downarrow$
  \item In-group cohesion strengthens $\Rightarrow$ $\Psi_S \uparrow$ locally
\end{itemize}

The coherence index $K$ changes differently across regions:

\begin{center}
\renewcommand{\arraystretch}{1.2}
\begin{tabular}{lcc}
\hline
\textbf{Region/Group} & \textbf{$\Delta K$} & \textbf{Mechanism} \\
\hline
Manufacturing heartland (USA) & $+$ & Policy matches identity \\
Coastal/export regions (USA) & $-$ & Policy conflicts with globalist identity \\
Nationalist base (China) & $+$ & External threat unifies \\
Export sector (China) & $-$ & Business model disrupted \\
WTO/global institutions & $-$ & Rules violated \\
\hline
\end{tabular}
\end{center}

\paragraph{Phase 3: Context Shift and Hysteresis ($t = 2$--$3$).}

The crucial prediction concerns $\Psi$-dynamics. Tariffs don't just change prices;
they shift context:

\begin{equation}
\Psi_{t+1} = \Psi_t + \eta \cdot (\text{tariff actions}_t) + \varepsilon_t
\label{eq:psi-dynamics-trade}
\end{equation}

Specifically:
\begin{itemize}
  \item $\Psi_{norm}$: ``Free trade is good'' norm weakens ($\eta_{norm} < 0$)
  \item $\Psi_I$: WTO institutional authority erodes ($\eta_I < 0$)
  \item $\Psi_M$: Supply chains permanently reconfigure ($\eta_M < 0$)
  \item $\Psi_S$: Nationalist sentiment strengthens within countries ($\eta_S > 0$ locally)
\end{itemize}

\subsubsection{The Hysteresis Prediction}

\begin{tcolorbox}[colback=blue!5!white, colframe=blue!75!black, title=Prediction 1: Trade War Hysteresis]
Trade wars have persistent effects beyond their direct economic costs because 
they shift $\Psi$. Even if tariffs are removed, the context has changed: 
trust is lower, supply chains have restructured, nationalist identities 
have strengthened. The global trading system moves to a different $C^*$---possibly permanently.

\textbf{Formal statement:}
\begin{equation}
\lim_{t \to \infty} \Psi_t \neq \Psi^{pre-war} \quad \text{even if tariffs removed at } t = T
\end{equation}

\textbf{Quantitative prediction:} A 30\% reduction in $\Psi_M$ from tariffs 
produces social welfare gain of $0.054$ SD and ecological welfare gain of 
$0.063$ SD, partially offsetting financial welfare loss of $0.084$ SD.
Net welfare effect depends on $\omega$ weights across dimensions.
\end{tcolorbox}

\subsubsection{Empirical Test Design}

\paragraph{Data Requirements.}
\begin{itemize}
  \item Pre-tariff $\Psi$ measures (2017): WTO compliance, trade openness, supply chain density
  \item Post-tariff $\Psi$ measures (2020, 2025): Same indicators
  \item FEPSDE outcomes by region: Income, life satisfaction, health, social cohesion, digital access
  \item Control: Regions with similar pre-tariff characteristics but different tariff exposure
\end{itemize}

\paragraph{Identification Strategy.}
Use variation in tariff exposure across industries and regions:
\begin{equation}
\Delta U_{d,r} = \alpha + \beta \cdot \text{TariffExposure}_r + \gamma \cdot \Delta \Psi_{r} + \varepsilon_{r}
\end{equation}

The framework predicts $\beta < 0$ for $U^F$, $\beta > 0$ for $U^S$.

\paragraph{Falsification Criteria.}
The prediction is falsified if:
\begin{enumerate}
  \item Tariff removal fully restores pre-war $\Psi$ within 2 years
  \item Social welfare ($U^S$) declines in tariff-imposing regions
  \item No persistent supply chain restructuring observed
\end{enumerate}

\subsubsection{Policy Implications}

\begin{enumerate}
  \item \textbf{Trade negotiations:} Cannot assume return to status quo ante---new equilibrium must be negotiated
  \item \textbf{Compensation design:} Financial compensation insufficient if identity and social gains are primary
  \item \textbf{Timing:} Early intervention prevents $\Psi$-lock-in; late intervention faces hysteresis
  \item \textbf{Regional policy:} Different regions experience different $K$ and $Q$ effects; uniform policy inappropriate
\end{enumerate}

%%%%%%%%%%%%%%%%%%%%%%%%%%%%%%%%%%%%%%%%%%%%%%%%%%%%%%%%%%%%%%%%%%%%%%%%%%%%%%%
\subsection{Prediction 2: Cross-Domain Coherence Spillovers}
\label{subsec:pred-spillovers}
%%%%%%%%%%%%%%%%%%%%%%%%%%%%%%%%%%%%%%%%%%%%%%%%%%%%%%%%%%%%%%%%%%%%%%%%%%%%%%%

\subsubsection{The Phenomenon}

Financial crises seem to cause political crises with a lag. The 2008 financial 
crisis was followed by Occupy Wall Street (2011), the Tea Party movement, 
Brexit (2016), and the Trump election (2016). Is this coincidence or mechanism?

\subsubsection{Standard Explanations}

\paragraph{Economic voting.} Financial distress causes voters to punish incumbents. 
But this predicts \emph{electoral change}, not \emph{institutional destabilization}.

\paragraph{Inequality.} Financial crises increase inequality, fueling resentment.
But inequality was rising before 2008; the crisis accelerated political change disproportionately.

\paragraph{Special interests.} Bailouts create moral hazard and public anger.
But anger alone doesn't explain the \emph{pattern} of cross-domain spillovers.

\subsubsection{Framework Mechanism: The $C$-Matrix Structure}

The complementarity matrix $C$ has cross-domain terms:

\begin{equation}
C = \begin{pmatrix}
C^{FF} & C^{FS} & C^{FP} & C^{FI} \\
C^{SF} & C^{SS} & C^{SP} & C^{SI} \\
C^{PF} & C^{PS} & C^{PP} & C^{PI} \\
C^{IF} & C^{IS} & C^{IP} & C^{II}
\end{pmatrix}
\end{equation}

where F = financial, S = social, P = political, I = institutional.

When $C^{FS} \neq 0$, financial instability directly affects social trust.
When $C^{SP} \neq 0$, social distrust feeds political instability.
The result is a \emph{cascade}:

\begin{equation}
\Delta K^F_t < 0 \xrightarrow{C^{FS}} \Delta K^S_{t+\tau} < 0 \xrightarrow{C^{SP}} \Delta K^P_{t+2\tau} < 0 \xrightarrow{C^{PI}} \Delta K^I_{t+3\tau} < 0
\end{equation}

\subsubsection{Quantitative Cascade Model}

Using the $\Gamma$-matrix, we can quantify the cascade. Financial crisis 
reduces $\Psi_E$ (economic stability) and $\Psi_I$ (institutional trust):

\paragraph{Stage 1: Financial $\to$ Social.}
\begin{equation}
\Delta U^S = \gamma_{S,E} \cdot \Delta \Psi_E + \gamma_{S,I} \cdot \Delta \Psi_I = 0.01 \cdot (-0.5) + 0.11 \cdot (-0.3) = -0.038
\end{equation}

\paragraph{Stage 2: Social $\to$ Emotional.}
The decline in social welfare reduces social trust ($\Psi_S$), which then affects emotional welfare:
\begin{equation}
\Delta U^E = \gamma_{E,S} \cdot \Delta \Psi_S = 0.47 \cdot (-0.2) = -0.094
\end{equation}

\paragraph{Stage 3: Emotional $\to$ Political.}
Reduced emotional wellbeing manifests as anger, resentment, and political alienation.
This creates demand for anti-establishment politics.

\paragraph{Cascade Multiplier.}
The total effect exceeds the sum of direct effects due to feedback:
\begin{equation}
\text{Total Effect} = \text{Direct} + \text{Cascade} = \Delta U^F + \sum_{d \neq F} \sum_{j} \gamma_{dj} \cdot \Delta \Psi_j(t+\tau)
\end{equation}

\subsubsection{The Prediction}

\begin{tcolorbox}[colback=blue!5!white, colframe=blue!75!black, title=Prediction 2: Cross-Domain Coherence Spillovers]
A negative shock to coherence in one domain predicts subsequent decline 
in coherence in other domains, even without direct causal links, with 
predictable lags and magnitudes.

\textbf{Specific prediction:} Financial crises ($\Delta K^F < 0$) cause political 
crises ($\Delta K^P < 0$) with a 5--8 year lag. The lag reflects the time 
required for:
\begin{enumerate}
  \item Financial $\to$ Social trust erosion (1--2 years)
  \item Social $\to$ Political alienation (2--3 years)
  \item Political $\to$ Electoral/institutional change (2--3 years)
\end{enumerate}

\textbf{Magnitude:} A 1 SD financial crisis shock produces approximately 
0.4 SD political instability effect after full cascade.
\end{tcolorbox}

\subsubsection{Historical Validation}

\begin{center}
\renewcommand{\arraystretch}{1.2}
\begin{tabular}{lccccc}
\hline
\textbf{Financial Crisis} & \textbf{Year} & \textbf{Social Effect} & \textbf{Political Effect} & \textbf{Lag} \\
\hline
Great Depression & 1929 & Labor unrest 1930s & Fascism/New Deal & 4--10 yrs \\
Latin America debt & 1982 & Social unrest & Democratic transitions & 5--8 yrs \\
Asian financial & 1997 & Social instability & Regime changes & 2--5 yrs \\
Global financial & 2008 & Occupy/Tea Party & Brexit/Trump & 8 yrs \\
\hline
\end{tabular}
\end{center}

The pattern is consistent: financial crises precede political transformations 
with multi-year lags, mediated by social trust erosion.

\subsubsection{Empirical Test Design}

\paragraph{Data.}
\begin{itemize}
  \item Financial coherence: Banking crisis indicator, credit spreads, volatility indices
  \item Social coherence: Trust surveys (GSS, Eurobarometer), social capital measures
  \item Political coherence: Vote fragmentation, institutional trust, protest frequency
  \item Panel: 50+ countries, 1970--2025
\end{itemize}

\paragraph{Model.}
Vector autoregression with coherence indices:
\begin{equation}
\mathbf{K}_t = A_1 \mathbf{K}_{t-1} + A_2 \mathbf{K}_{t-2} + \ldots + \varepsilon_t
\end{equation}

The framework predicts off-diagonal elements of $A$ are significant, 
with financial $\to$ social $\to$ political ordering in impulse responses.

\paragraph{Falsification.}
The prediction is falsified if:
\begin{enumerate}
  \item Financial crises show no predictive power for subsequent political instability
  \item The lag structure is random rather than systematic
  \item Social trust does not mediate the financial-political relationship
\end{enumerate}

\subsubsection{Policy Implications}

\begin{enumerate}
  \item \textbf{Financial regulation:} Preventing financial crises prevents political crises
  \item \textbf{Crisis response:} Social trust maintenance as important as financial stabilization
  \item \textbf{Political forecasting:} Financial crisis today predicts political crisis in 5--8 years
  \item \textbf{Institutional design:} Firebreaks between domains reduce cascade risk
\end{enumerate}

%%%%%%%%%%%%%%%%%%%%%%%%%%%%%%%%%%%%%%%%%%%%%%%%%%%%%%%%%%%%%%%%%%%%%%%%%%%%%%%
\subsection{Prediction 3: Asymmetric Recovery and the Coherence Restoration Function}
\label{subsec:pred-asymmetric}
%%%%%%%%%%%%%%%%%%%%%%%%%%%%%%%%%%%%%%%%%%%%%%%%%%%%%%%%%%%%%%%%%%%%%%%%%%%%%%%

\subsubsection{The Puzzle}

Some crises resolve quickly; others persist for decades. Japan's ``lost decades'' 
contrast with rapid recoveries elsewhere. Why do similar-magnitude shocks 
produce different recovery trajectories?

\subsubsection{Standard Explanations}

\paragraph{Policy mistakes.} Japan's slow recovery reflects monetary policy errors.
But similar policies produced different outcomes elsewhere.

\paragraph{Structural factors.} Demographics, debt levels, etc.
But these don't explain the \emph{non-linearity} of recovery.

\subsubsection{Framework Mechanism: Convex Recovery Costs}

Define the coherence restoration function:
\begin{equation}
\tau(K) = \text{time to restore } K \text{ to } K^* = 1
\end{equation}

The framework predicts $\tau(K)$ is \emph{convex}:
\begin{equation}
\tau''(K) > 0 \quad \Rightarrow \quad \text{Larger coherence gaps take disproportionately longer to close}
\end{equation}

\paragraph{Mechanism 1: Complementarity Breakdown.}
When $K$ falls below threshold $\underline{K}$, complementarities that 
supported the old equilibrium break down. The system must build \emph{new} 
complementarities, not restore old ones.

\paragraph{Mechanism 2: Belief Coordination.}
Recovery requires coordinated belief revision. With many agents, coordination 
becomes harder as beliefs diverge more---a convex relationship.

\paragraph{Mechanism 3: Institutional Erosion.}
Low $K$ erodes institutions ($\Psi_I$ falls). Institutional rebuilding is 
slow and path-dependent, creating hysteresis in recovery.

\subsubsection{Quantitative Model}

The recovery dynamics follow:
\begin{equation}
\frac{dK}{dt} = \phi \cdot (K^* - K) \cdot \Psi_I(K) \cdot \Psi_S(K)
\label{eq:recovery-dynamics}
\end{equation}

where $\Psi_I(K)$ and $\Psi_S(K)$ are themselves functions of coherence.

Using the $\Gamma$-matrix:
\begin{itemize}
  \item $\Psi_I$ depends on financial welfare: $\Psi_I = f(U^F) = f(\gamma_{F,I} \cdot \Psi_I + \ldots)$
  \item $\Psi_S$ depends on social welfare: $\Psi_S = g(U^S) = g(\gamma_{S,S} \cdot \Psi_S + \ldots)$
\end{itemize}

This creates feedback loops that slow recovery when $K$ is low.

\subsubsection{The Prediction}

\begin{tcolorbox}[colback=blue!5!white, colframe=blue!75!black, title=Prediction 3: Asymmetric Recovery]
Recovery from coherence shocks is asymmetric: small shocks recover quickly, 
large shocks recover slowly and may not recover at all.

\textbf{Formal statement:}
\begin{equation}
\frac{\tau(K_1)}{\tau(K_2)} > \frac{K^* - K_1}{K^* - K_2} \quad \text{for } K_1 < K_2
\end{equation}

\textbf{Threshold effect:} Below critical coherence $\underline{K} \approx 0.6$, 
recovery requires ``big push'' coordinated intervention. Marginal reforms 
are insufficient.

\textbf{Quantitative prediction:} A crisis reducing $K$ by 0.2 (mild) requires 
$\sim$3 years to recover. A crisis reducing $K$ by 0.4 (severe) requires 
$\sim$15 years---not 6 years as linear extrapolation would suggest.
\end{tcolorbox}

\subsubsection{Case Study: Japan vs. Nordic Countries}

\paragraph{Japan 1990--2020.}
\begin{itemize}
  \item Initial shock: Asset bubble collapse, $\Delta K \approx -0.4$
  \item $\Psi_I$ erosion: Banking system paralysis, policy gridlock
  \item $\Psi_S$ erosion: Social contract fraying, declining trust
  \item Recovery: Incomplete after 30 years
\end{itemize}

\paragraph{Nordic Countries 1990--2000.}
\begin{itemize}
  \item Initial shock: Banking crisis, $\Delta K \approx -0.25$
  \item $\Psi_I$ maintained: Swift policy response, institutional reform
  \item $\Psi_S$ maintained: Social compact preserved
  \item Recovery: Complete within 5--7 years
\end{itemize}

The difference: Japan crossed the $\underline{K}$ threshold; Nordic countries did not.

\subsubsection{Empirical Test Design}

\paragraph{Data.}
\begin{itemize}
  \item Crisis magnitude: GDP decline, financial sector losses, unemployment spike
  \item Recovery time: Years to return to pre-crisis trend
  \item Coherence proxies: Institutional quality indices, social trust, policy consensus
\end{itemize}

\paragraph{Model.}
Non-linear regression:
\begin{equation}
\log(\tau_i) = \alpha + \beta_1 \cdot \text{CrisisMagnitude}_i + \beta_2 \cdot \text{CrisisMagnitude}_i^2 + \gamma X_i + \varepsilon_i
\end{equation}

The framework predicts $\beta_2 > 0$ (convexity).

\paragraph{Falsification.}
The prediction is falsified if:
\begin{enumerate}
  \item Recovery time is linear in crisis magnitude
  \item No threshold effects are observed
  \item Institutional quality does not mediate recovery speed
\end{enumerate}

\subsubsection{Policy Implications}

\begin{enumerate}
  \item \textbf{Crisis response:} Speed matters more than magnitude---prevent crossing $\underline{K}$
  \item \textbf{Intervention size:} Small crises need proportional response; large crises need disproportionately large response
  \item \textbf{Institutional preservation:} Maintaining $\Psi_I$ and $\Psi_S$ during crisis accelerates recovery
  \item \textbf{Coordination:} Below $\underline{K}$, decentralized adjustment fails; coordinated ``big push'' required
\end{enumerate}

%%%%%%%%%%%%%%%%%%%%%%%%%%%%%%%%%%%%%%%%%%%%%%%%%%%%%%%%%%%%%%%%%%%%%%%%%%%%%%%
\subsection{Prediction 4: Technological Disruption and the Techlash}
\label{subsec:pred-techlash}
%%%%%%%%%%%%%%%%%%%%%%%%%%%%%%%%%%%%%%%%%%%%%%%%%%%%%%%%%%%%%%%%%%%%%%%%%%%%%%%

\subsubsection{The Phenomenon}

Technological progress is supposed to increase welfare. Yet recent years 
have seen ``techlash''---public backlash against technology companies, 
concerns about automation, and nostalgia for pre-digital life. Why would 
progress produce discontent?

\subsubsection{Standard Explanations}

\paragraph{Luddite fallacy.} People irrationally fear technology; they'll adapt.
But current concerns persist despite adaptation.

\paragraph{Transition costs.} Short-run disruption, long-run gains.
But timeline for ``long run'' keeps extending.

\paragraph{Inequality.} Technology benefits accrue to owners.
True, but doesn't explain the \emph{identity} dimension of backlash.

\subsubsection{Framework Mechanism: $\Psi$-Disruption}

Technological change ($\Delta \Psi_T$ where T = technology) affects all 
FEPSDE dimensions through the $\Gamma$-matrix. But crucially, technology 
also shifts \emph{other} $\Psi$ dimensions:

\begin{equation}
\Psi_{t+1} = \Psi_t + \Phi \cdot \Delta \Psi^{tech}_t
\end{equation}

where $\Phi$ is the technology-context transmission matrix:

\begin{center}
\renewcommand{\arraystretch}{1.2}
\begin{tabular}{lcc}
\hline
\textbf{Context Dimension} & \textbf{$\Phi$ Effect} & \textbf{Mechanism} \\
\hline
$\Psi_K$ (Information) & $+$ & More data, more access \\
$\Psi_M$ (Market scope) & $+$ & Global platforms \\
$\Psi_F$ (Flexibility) & $+/-$ & Gig economy: flexible but precarious \\
$\Psi_S$ (Social trust) & $-$ & Filter bubbles, atomization \\
$\Psi_T$ (Temporal stability) & $-$ & Accelerating change, uncertainty \\
\hline
\end{tabular}
\end{center}

\paragraph{Net Welfare Effect.}
Using the $\Gamma$-matrix:
\begin{equation}
\Delta Q = \sum_{d} \sum_{j} \omega_d \cdot \gamma_{dj} \cdot \Delta \Psi_j
\end{equation}

Technology increases $\Psi_K$ and $\Psi_M$ (positive for $U^F$, $U^D$) 
but decreases $\Psi_S$ and $\Psi_T$ (negative for $U^E$, $U^S$).

\paragraph{Numerical Example: Digital Platform Expansion.}

Assume digital platforms produce:
\begin{itemize}
  \item $\Delta \Psi_K = +0.3$ (more information)
  \item $\Delta \Psi_M = +0.2$ (market expansion)
  \item $\Delta \Psi_S = -0.2$ (social atomization)
  \item $\Delta \Psi_T = -0.15$ (faster change)
\end{itemize}

Effects by dimension:
\begin{align*}
\Delta U^F &= 0.25 \cdot (+0.3) + 0.28 \cdot (+0.2) = +0.131 \\
\Delta U^E &= 0.09 \cdot (+0.3) + 0.47 \cdot (-0.2) + 0.29 \cdot (-0.15) = -0.110 \\
\Delta U^S &= 0.08 \cdot (+0.3) + 0.51 \cdot (-0.2) + (-0.18) \cdot (+0.2) = -0.114 \\
\Delta U^D &= 0.39 \cdot (+0.3) + 0.24 \cdot (+0.2) = +0.165
\end{align*}

\textbf{Result:} Financial and digital welfare increase; emotional and social welfare decrease.
Net effect depends on weights $\omega_d$---and these weights vary across populations.

\subsubsection{The Prediction}

\begin{tcolorbox}[colback=blue!5!white, colframe=blue!75!black, title=Prediction 4: Techlash as Coherence Response]
Rapid technological change triggers ``techlash'' not despite its benefits 
but \emph{because} of asymmetric effects across welfare dimensions. Groups 
whose identity is tied to pre-digital social structures ($U^S$, $U^{IDN}$) 
experience net welfare loss even as aggregate $U^F$ rises.

\textbf{Quantitative prediction:} Populations with $\omega_S > 0.25$ 
(high social welfare weight) will exhibit net negative response to 
digitalization, while populations with $\omega_F > 0.35$ (high financial 
weight) will exhibit positive response.

\textbf{Demographic prediction:} Techlash will be concentrated among:
\begin{itemize}
  \item Older populations (higher $\omega_S$, identity tied to pre-digital)
  \item Rural populations (stronger community ties, lower digital gains)
  \item Workers in disrupted industries (identity threat + financial loss)
\end{itemize}
\end{tcolorbox}

\subsubsection{The Coherence Dimension}

Technology doesn't just affect welfare levels; it disrupts coherence.
When technology changes faster than beliefs and institutions can adapt:

\begin{equation}
\frac{d\Psi^{tech}}{dt} > \frac{dC^*}{dt} \quad \Rightarrow \quad K \downarrow
\end{equation}

The system is in disequilibrium: old complementarities (skills, institutions, 
social networks) no longer fit the new technological context.

\paragraph{Examples.}
\begin{itemize}
  \item Labor markets: Skills mismatch $\Rightarrow$ $K^{labor} \downarrow$
  \item Journalism: Business model collapse $\Rightarrow$ $K^{media} \downarrow$
  \item Retail: Physical stores vs. e-commerce $\Rightarrow$ $K^{retail} \downarrow$
  \item Democracy: Social media vs. legacy institutions $\Rightarrow$ $K^{political} \downarrow$
\end{itemize}

\subsubsection{Empirical Test Design}

\paragraph{Data.}
\begin{itemize}
  \item Technology adoption: Broadband penetration, smartphone usage, platform engagement
  \item FEPSDE outcomes: Income, life satisfaction, health, social connection, digital participation
  \item Attitudes: Technology optimism/pessimism, support for regulation
  \item Demographics: Age, location, occupation, education
\end{itemize}

\paragraph{Model.}
Heterogeneous effects regression:
\begin{equation}
\text{Techlash}_i = \alpha + \beta \cdot \text{TechExposure}_i + \gamma \cdot (\text{TechExposure}_i \times \omega_{S,i}) + \delta X_i + \varepsilon_i
\end{equation}

The framework predicts $\gamma > 0$: higher social welfare weight increases techlash.

\paragraph{Falsification.}
The prediction is falsified if:
\begin{enumerate}
  \item Techlash is uniform across demographic groups
  \item Social welfare weight does not predict technology attitudes
  \item Technology adoption improves $U^S$ and $U^E$ as well as $U^F$ and $U^D$
\end{enumerate}

\subsubsection{Policy Implications}

\begin{enumerate}
  \item \textbf{Technology policy:} Cannot ignore non-financial welfare dimensions
  \item \textbf{Transition support:} Address identity and social effects, not just retraining
  \item \textbf{Pace management:} Controlled adoption may dominate rapid disruption
  \item \textbf{Platform design:} Build in social coherence preservation, not just engagement optimization
\end{enumerate}

%%%%%%%%%%%%%%%%%%%%%%%%%%%%%%%%%%%%%%%%%%%%%%%%%%%%%%%%%%%%%%%%%%%%%%%%%%%%%%%
\subsection{Prediction 5: Identity Cannot Be Compensated Financially}
\label{subsec:pred-identity}
%%%%%%%%%%%%%%%%%%%%%%%%%%%%%%%%%%%%%%%%%%%%%%%%%%%%%%%%%%%%%%%%%%%%%%%%%%%%%%%

\subsubsection{The Policy Puzzle}

Coal miners offered retraining programs often don't take them. Displaced 
manufacturing workers reject service sector jobs at higher wages. 
``Economic anxiety'' voters don't respond to economic programs. Why?

\subsubsection{Standard Explanations}

\paragraph{Information failure.} Workers don't know about opportunities.
But information programs show limited effectiveness.

\paragraph{Moving costs.} Geographic/skill barriers.
But subsidies to cover costs don't increase take-up proportionally.

\paragraph{Irrationality.} Nostalgia, stubbornness.
This is description, not explanation.

\subsubsection{Framework Mechanism: Identity Utility}

The framework adds $U^{IDN}$ as a distinct welfare dimension (Chapter~8).
Identity utility depends on alignment between behavior and self-concept:

\begin{equation}
U^{IDN} = \omega_{IDN} \cdot \text{Match}(a, I)
\end{equation}

where $a$ is action and $I$ is identity.

\paragraph{Non-Fungibility.}
The critical insight: identity utility is \emph{not fungible} with financial utility.
Offering higher $U^F$ cannot compensate for lower $U^{IDN}$ if the identity 
dimension is separately weighted:

\begin{equation}
Q = \omega_F \cdot U^F + \omega_{IDN} \cdot U^{IDN} + \ldots
\end{equation}

A coal miner's identity is ``miner,'' not ``worker in energy sector.'' 
Retraining to solar panel installation offers comparable $U^F$ but 
threatens $U^{IDN}$.

\subsubsection{Quantifying the Identity Premium}

Using revealed preference from job choice data, we can estimate the 
identity premium: the financial sacrifice people accept to preserve 
identity-consistent employment.

\paragraph{Model.}
Agent chooses between:
\begin{itemize}
  \item Option A: Identity-consistent job, wage $w_A$
  \item Option B: Identity-inconsistent job, wage $w_B > w_A$
\end{itemize}

Choice of A reveals:
\begin{equation}
\omega_F \cdot w_A + \omega_{IDN} \cdot U^{IDN}_A > \omega_F \cdot w_B + \omega_{IDN} \cdot U^{IDN}_B
\end{equation}

Solving:
\begin{equation}
\frac{\omega_{IDN}}{\omega_F} > \frac{w_B - w_A}{U^{IDN}_A - U^{IDN}_B}
\end{equation}

\paragraph{Empirical Estimates.}
Studies suggest identity premiums of 15--40\% of wages for occupation-specific 
identity \citep{PAP-akerlof2000identity}. This implies:
\begin{equation}
\frac{\omega_{IDN}}{\omega_F} \approx 0.2 \text{ to } 0.5
\end{equation}

\subsubsection{The Prediction}

\begin{tcolorbox}[colback=blue!5!white, colframe=blue!75!black, title=Prediction 5: Identity Non-Fungibility]
Economic compensation programs that address only $U^F$ will fail for 
populations where $U^{IDN}$ is the binding constraint. Financial transfers 
cannot substitute for identity preservation.

\textbf{Quantitative prediction:} Retraining program take-up will be:
\begin{itemize}
  \item High ($>60\%$) when new occupation preserves identity category
  \item Low ($<30\%$) when new occupation violates identity category
  \item Intermediate when identity implications are ambiguous
\end{itemize}

\textbf{Compensation threshold:} Identity-inconsistent transitions require 
wage premiums of $>30\%$ to achieve equivalent take-up to identity-consistent 
transitions with no premium.
\end{tcolorbox}

\subsubsection{Case Studies}

\paragraph{Coal-to-Solar Transition.}
\begin{itemize}
  \item Identity category: ``Miner'' $\neq$ ``Solar technician''
  \item Financial offer: Comparable wages, retraining subsidies
  \item Take-up: Low despite financial equivalence
  \item Framework explanation: $\Delta U^{IDN} < 0$ dominates $\Delta U^F \approx 0$
\end{itemize}

\paragraph{Manufacturing-to-Healthcare.}
\begin{itemize}
  \item Identity category: ``Maker'' $\neq$ ``Caregiver'' (gendered dimension)
  \item Financial offer: Healthcare often pays more
  \item Take-up: Very low for male manufacturing workers
  \item Framework explanation: Gender identity compounds occupational identity loss
\end{itemize}

\paragraph{Successful Transition: German Reunification.}
\begin{itemize}
  \item Approach: Maintained industry identity where possible (e.g., Optics in Jena)
  \item Key: Framed transition as ``modernization'' not ``abandonment''
  \item Result: Higher take-up than pure retraining programs
  \item Framework explanation: $U^{IDN}$ partially preserved
\end{itemize}

\subsubsection{Empirical Test Design}

\paragraph{Data.}
\begin{itemize}
  \item Displaced worker surveys: Occupation, offered alternatives, choices made
  \item Identity measures: Occupational identity scale, community attachment
  \item Wage data: Offered wages in current vs. alternative occupations
  \item Program outcomes: Take-up rates, completion rates, employment outcomes
\end{itemize}

\paragraph{Model.}
Structural choice model with identity:
\begin{equation}
P(\text{Accept}_i) = \Phi\left(\beta_F \cdot \Delta w_i + \beta_{IDN} \cdot \Delta \text{Identity}_i + \gamma X_i\right)
\end{equation}

The framework predicts $\beta_{IDN}$ is significant and $\beta_{IDN}/\beta_F \approx 0.2$--$0.5$.

\paragraph{Falsification.}
The prediction is falsified if:
\begin{enumerate}
  \item Wage differences fully explain transition choices
  \item Identity-consistent vs. inconsistent transitions show no take-up difference
  \item Higher financial compensation eliminates identity effects
\end{enumerate}

\subsubsection{Policy Implications}

\begin{enumerate}
  \item \textbf{Program design:} Incorporate identity continuity, not just skills and wages
  \item \textbf{Framing:} ``Modernization within tradition'' beats ``complete transformation''
  \item \textbf{Community approach:} Collective transitions preserve social identity better than individual
  \item \textbf{Gradual transitions:} Bridge occupations that maintain partial identity
  \item \textbf{Recognition:} Honor past identity while building new one
\end{enumerate}

%%%%%%%%%%%%%%%%%%%%%%%%%%%%%%%%%%%%%%%%%%%%%%%%%%%%%%%%%%%%%%%%%%%%%%%%%%%%%%%
\subsection{Prediction 6: The Coherence Trap and Development}
\label{subsec:pred-coherence-trap}
%%%%%%%%%%%%%%%%%%%%%%%%%%%%%%%%%%%%%%%%%%%%%%%%%%%%%%%%%%%%%%%%%%%%%%%%%%%%%%%

\subsubsection{The Development Puzzle}

Many countries reach middle-income status, then stagnate. The ``middle-income 
trap'' has caught countries from Latin America to Southeast Asia. Why is 
the transition from middle to high income so difficult?

\subsubsection{Standard Explanations}

\paragraph{Institutional quality.} Weak institutions prevent sustained growth.
But institutions that enabled middle-income growth now fail---why the discontinuity?

\paragraph{Human capital.} Education systems don't produce innovation.
But education quality doesn't change discontinuously at middle income.

\paragraph{Industrial policy.} Export manufacturing hits limits.
But this describes the trap, not explains why escape is hard.

\subsubsection{Framework Mechanism: High-K, Low-Q Equilibria}

The framework distinguishes coherence ($K$) from quality ($Q$).
A system can be highly coherent around a \emph{bad} equilibrium:

\begin{center}
\renewcommand{\arraystretch}{1.3}
\begin{tabular}{lcc}
\hline
\textbf{Configuration} & \textbf{$K$} & \textbf{$Q$} \\
\hline
Middle-income trap & High & Low-Medium \\
High-income equilibrium & High & High \\
Transition & \textbf{Low} & Variable \\
\hline
\end{tabular}
\end{center}

\paragraph{The Middle-Income Configuration.}
The middle-income trap is a coherent equilibrium:
\begin{itemize}
  \item Export manufacturing fits low wages
  \item Low wages fit weak domestic demand
  \item Weak demand fits export orientation
  \item Authoritarian stability fits rapid industrialization
\end{itemize}

All elements reinforce each other: $K \approx 1$.

\paragraph{The Transition Problem.}
Moving to high income requires:
\begin{itemize}
  \item Innovation (not just manufacturing)
  \item Higher wages (strong domestic demand)
  \item Institutional quality (property rights, rule of law)
  \item Political openness (creative destruction requires tolerance of disruption)
\end{itemize}

But these changes are \emph{complementary}---each requires the others.
Changing one without the others creates incoherence:

\begin{equation}
C^{high-income} = \begin{pmatrix}
\text{Innovation} & + & + & + \\
+ & \text{High wages} & + & + \\
+ & + & \text{Institutions} & + \\
+ & + & + & \text{Openness}
\end{pmatrix}
\end{equation}

Partial reform breaks current coherence without establishing new coherence:
$K_{transition} < K_{trap}$.

\subsubsection{The Valley of Incoherence}

\begin{equation}
Q_{transition}(t) = Q_{trap} + \underbrace{(\Delta Q_{potential})}_{\text{long-run gain}} - \underbrace{(1 - K(t)) \cdot \text{IncoherenceCost}}_{\text{transition cost}}
\end{equation}

During transition:
\begin{itemize}
  \item Current welfare falls (incoherence costs)
  \item Uncertainty is high (which equilibrium will emerge?)
  \item Political opposition: ``Things were better before''
  \item Temptation to retreat to old equilibrium
\end{itemize}

\subsubsection{Quantifying the Trap Using $\Gamma$-Matrix}

The $\Gamma$-matrix shows how context affects welfare. The middle-income 
trap is characterized by:
\begin{itemize}
  \item Moderate $\Psi_E$ (economic development)
  \item Low $\Psi_I$ (institutional quality)
  \item Low $\Psi_C$ (cognitive/innovation capacity)
  \item Moderate $\Psi_F$ (factor flexibility)
\end{itemize}

From the $\Gamma$-matrix:
\begin{align*}
U^F &= 0.31 \cdot \Psi_I + 0.42 \cdot \Psi_E + 0.33 \cdot \Psi_F + \ldots \\
&= 0.31 \cdot (0.4) + 0.42 \cdot (0.5) + 0.33 \cdot (0.5) + \ldots \\
&\approx \text{Medium financial welfare (trapped)}
\end{align*}

To reach high income requires increasing $\Psi_I$ and $\Psi_C$---but these 
are the dimensions most complementary with each other.

\subsubsection{The Prediction}

\begin{tcolorbox}[colback=blue!5!white, colframe=blue!75!black, title=Prediction 6: The Coherence Trap]
Development is not continuous but involves discrete jumps between coherent 
configurations. The ``middle-income trap'' is a coherence trap: the current 
equilibrium is stable ($K \approx 1$) but suboptimal ($Q < Q^{max}$).

\textbf{Escape requires:} Coordinated, simultaneous change across multiple 
complementary dimensions---a ``big push''---not gradual reform.

\textbf{Specific prediction:} Countries attempting gradual reform will show:
\begin{enumerate}
  \item Initial improvement followed by stagnation or reversal
  \item Reform reversal probability $>50\%$ within 10 years
  \item Successful escapes characterized by rapid, multi-dimensional reform
\end{enumerate}

\textbf{Threshold:} Escape requires simultaneous improvement in at least 
3 of 4 dimensions: institutions ($\Psi_I$), innovation ($\Psi_C$), 
openness ($\Psi_S$), and market flexibility ($\Psi_F$).
\end{tcolorbox}

\subsubsection{Historical Cases}

\begin{center}
\renewcommand{\arraystretch}{1.2}
\begin{tabular}{lcccc}
\hline
\textbf{Country} & \textbf{Approach} & \textbf{Dimensions Changed} & \textbf{Outcome} \\
\hline
South Korea & Big push (1960s--80s) & 4/4 & Escaped \\
Taiwan & Big push (1960s--80s) & 4/4 & Escaped \\
Brazil & Gradual reform & 1--2/4 & Trapped \\
Argentina & Partial reform, reversal & 2/4 $\to$ 1/4 & Trapped \\
Malaysia & Moderate reform & 2--3/4 & Upper-middle \\
Thailand & Moderate reform & 2/4 & Trapped \\
\hline
\end{tabular}
\end{center}

\subsubsection{Empirical Test Design}

\paragraph{Data.}
\begin{itemize}
  \item Income level trajectories: GDP per capita, 1960--2025
  \item $\Psi$ dimensions: Institutional quality, education/innovation, trade openness, labor flexibility
  \item Reform events: Policy changes, structural adjustments
  \item Coherence proxies: Policy consistency, institutional stability
\end{itemize}

\paragraph{Model.}
Threshold regression:
\begin{equation}
P(\text{Escape}_i) = \Phi\left(\beta \cdot \sum_j \mathbf{1}[\Delta \Psi_{j,i} > \bar{\Delta}]\right)
\end{equation}

The framework predicts non-linear relationship: escape probability increases 
sharply when 3+ dimensions show substantial improvement.

\paragraph{Falsification.}
The prediction is falsified if:
\begin{enumerate}
  \item Gradual reform is as effective as big push
  \item Single-dimension reform produces sustained escape
  \item No threshold effect in number of dimensions reformed
\end{enumerate}

\subsubsection{Policy Implications}

\begin{enumerate}
  \item \textbf{Comprehensive reform:} Single-issue reform fails; package required
  \item \textbf{Speed:} Rapid transition minimizes time in incoherence valley
  \item \textbf{Commitment:} Credible commitment prevents retreat to trap
  \item \textbf{External support:} International institutions can anchor reform commitment
  \item \textbf{Political economy:} Build coalition that benefits from new equilibrium
\end{enumerate}

%%%%%%%%%%%%%%%%%%%%%%%%%%%%%%%%%%%%%%%%%%%%%%%%%%%%%%%%%%%%%%%%%%%%%%%%%%%%%%%
\subsection{Prediction 7: Multi-Level Coherence Conflicts}
\label{subsec:pred-multilevel}
%%%%%%%%%%%%%%%%%%%%%%%%%%%%%%%%%%%%%%%%%%%%%%%%%%%%%%%%%%%%%%%%%%%%%%%%%%%%%%%

\subsubsection{The Climate Puzzle}

Every country agrees climate change is a problem. International agreements 
are signed. Yet emissions continue rising. Why does collective agreement 
fail to produce collective action?

\subsubsection{Standard Explanations}

\paragraph{Free-rider problem.} Each country benefits from others' abatement.
But this doesn't explain \emph{intensity} of disagreement and conflict.

\paragraph{Discounting.} Future costs vs. present benefits.
But countries agree on the problem's severity.

\paragraph{Sovereignty.} No global enforcement mechanism.
But this is institutional description, not causal explanation.

\subsubsection{Framework Mechanism: Coherence at Different Levels}

The framework allows different coherence levels at different scales:

\begin{align}
K^{national} &= \text{Coherence within country} \\
K^{global} &= \text{Coherence across countries}
\end{align}

\textbf{The key insight:} These can conflict. Optimal coherence at one level 
may destroy coherence at another.

\paragraph{Climate as Multi-Level Coherence Problem.}

\begin{center}
\renewcommand{\arraystretch}{1.3}
\begin{tabular}{lcc}
\hline
\textbf{Level} & \textbf{Configuration} & \textbf{$K$} \\
\hline
Individual country & Pursue national interest & High \\
Global system & Collective action on climate & Low \\
\hline
\end{tabular}
\end{center}

Each country has a coherent national strategy: economic growth, energy security, 
competitive advantage. These strategies are mutually consistent \emph{within} 
the country ($K^{national} \approx 1$).

But the \emph{global} system is incoherent: the sum of national strategies 
produces climate catastrophe ($K^{global} \ll 1$).

\subsubsection{The $\Gamma$-Matrix Perspective}

Climate action involves trade-offs visible in the $\Gamma$-matrix:

\begin{itemize}
  \item Climate action $\Rightarrow$ $\Delta \Psi_E < 0$ (economic costs)
  \item From $\Gamma$: $\gamma_{F,E} = 0.42$ (large positive)
  \item Therefore: $\Delta U^F = 0.42 \cdot \Delta \Psi_E < 0$
\end{itemize}

But climate damages affect ecological welfare:
\begin{itemize}
  \item No action $\Rightarrow$ $\Delta \Psi_{Eco} < 0$ (environmental degradation)
  \item From $\Gamma$: $\gamma_{Eco,Eco} = 1$ (direct)
  \item Therefore: $\Delta U^{Eco} < 0$
\end{itemize}

\paragraph{Level Asymmetry.}
Financial costs are \emph{national}; ecological costs are \emph{global}.
Each country optimizes:
\begin{equation}
Q^{national} = \omega_F \cdot U^F + \omega_{Eco}^{national} \cdot U^{Eco}_{national}
\end{equation}

But $U^{Eco}_{global}$ doesn't enter the national calculation---a coordination failure.

\subsubsection{The Prediction}

\begin{tcolorbox}[colback=blue!5!white, colframe=blue!75!black, title=Prediction 7: Multi-Level Coherence Conflict]
Policies that optimize coherence at one level may destroy coherence at 
another level. The climate crisis is fundamentally a multi-level coherence 
problem: high $K^{national}$, low $K^{global}$.

\textbf{General principle:}
\begin{equation}
\max_{policy} K^{level_1} \neq \max_{policy} K^{level_2} \quad \text{in general}
\end{equation}

\textbf{Specific predictions:}
\begin{enumerate}
  \item International agreements will be signed but not enforced (high $K^{rhetoric}$, low $K^{action}$)
  \item Countries with high $\omega_{Eco}$ will lead action; high $\omega_F$ will free-ride
  \item Solution requires either: (a) global governance ($K^{global}$ dominates), or (b) alignment of $Q^{national}$ with $Q^{global}$ (co-benefits)
\end{enumerate}

\textbf{Quantitative:} Climate cooperation will remain below optimal until 
at least one condition holds:
\begin{itemize}
  \item Global governance with enforcement ($K^{global}$ binding)
  \item Green technology cost reduction such that $\Delta U^F > 0$ (co-benefit alignment)
  \item Climate damages severe enough to shift $\omega_{Eco}^{national}$ upward
\end{itemize}
\end{tcolorbox}

\subsubsection{Other Multi-Level Conflicts}

\paragraph{Corporate vs. Regional Coherence.}
\begin{itemize}
  \item Company: Coherent strategy to offshore manufacturing $\Rightarrow K^{org} \uparrow$
  \item Region: Loses jobs, tax base, community $\Rightarrow K^{region} \downarrow$
  \item Result: Rational corporate decisions produce regional incoherence
\end{itemize}

\paragraph{Individual vs. Collective Rationality.}
\begin{itemize}
  \item Individual: Coherent to drive alone (convenient) $\Rightarrow K^{individual} = 1$
  \item City: Traffic congestion, pollution $\Rightarrow K^{city} \downarrow$
  \item Solution: Congestion pricing aligns individual and collective $K$
\end{itemize}

\subsubsection{Empirical Test Design}

\paragraph{Data.}
\begin{itemize}
  \item Climate commitments: Paris Agreement pledges, NDCs
  \item Actual emissions: National and global trajectories
  \item Policy coherence: Domestic policy vs. international commitments
  \item Level indicators: National vs. global welfare weights
\end{itemize}

\paragraph{Model.}
Multi-level coherence regression:
\begin{equation}
\text{EmissionsGap}_i = \alpha + \beta_1 \cdot K^{national}_i + \beta_2 \cdot K^{global} + \gamma \cdot (\omega_F/\omega_{Eco})_i + \varepsilon_i
\end{equation}

The framework predicts $\beta_1 < 0$ (higher national coherence $\Rightarrow$ larger gap) 
and $\gamma > 0$ (higher financial weight $\Rightarrow$ larger gap).

\paragraph{Falsification.}
The prediction is falsified if:
\begin{enumerate}
  \item High national coherence countries show better climate performance
  \item No systematic relationship between welfare weights and emissions
  \item International agreements produce proportional emissions reductions
\end{enumerate}

\subsubsection{Policy Implications}

\begin{enumerate}
  \item \textbf{Governance design:} Create institutions where $K^{global}$ is meaningful
  \item \textbf{Co-benefit strategies:} Align $Q^{national}$ with $Q^{global}$ through technology, trade
  \item \textbf{Compensation mechanisms:} Side payments to align national incentives
  \item \textbf{Moral framing:} Shift $\omega_{Eco}^{national}$ through norm entrepreneurship
  \item \textbf{Binding commitments:} Treaties with enforcement raise cost of defection
\end{enumerate}

%%%%%%%%%%%%%%%%%%%%%%%%%%%%%%%%%%%%%%%%%%%%%%%%%%%%%%%%%%%%%%%%%%%%%%%%%%%%%%%
\subsection{Prediction 8: Mechanism Design Must Model Context Dynamics}
\label{subsec:pred-mechanism}
%%%%%%%%%%%%%%%%%%%%%%%%%%%%%%%%%%%%%%%%%%%%%%%%%%%%%%%%%%%%%%%%%%%%%%%%%%%%%%%

\subsubsection{The Design Puzzle}

Market design has produced remarkable successes: kidney exchange, school 
choice, spectrum auctions. But some designed markets fail over time---
initial success followed by erosion. Why do well-designed mechanisms 
sometimes unravel?

\subsubsection{Standard Mechanism Design}

Classical mechanism design \citep{hurwicz1960, myerson1981, maskin1999} asks:
given agents with private information and strategic incentives, how can we 
design rules to achieve desired outcomes?

Key assumption: \textbf{Preferences are fixed.}

The mechanism takes preferences as given and optimizes rules.

\subsubsection{Framework Extension: Mechanisms Change Context}

The framework adds a crucial insight: \textbf{mechanisms change the context 
in which they operate.}

\begin{equation}
\Psi_{t+1} = f(\Psi_t, \text{mechanism}_t)
\end{equation}

A mechanism that works at $\Psi_t$ may not work at $\Psi_{t+1}$---because 
the mechanism itself changed $\Psi$.

\paragraph{Example: Trust Crowding.}
\begin{itemize}
  \item Context: Low-trust community, cooperation problem
  \item Classical solution: Introduce contracts and enforcement ($\Psi_I \uparrow$)
  \item Short run: Works! Cooperation increases
  \item Long run: People learn ``I don't need trust, I have contracts''
  \item Result: $\Psi_S \downarrow$ (social trust erodes)
  \item New problem: When contracts are incomplete or enforcement fails, cooperation collapses entirely
\end{itemize}

From the $\Gamma$-matrix: $\gamma_{S,I} = 0.11$ (institutional quality helps social welfare), 
but this is a \emph{static} relationship. The \emph{dynamic} effect of formal 
institutions substituting for informal trust may be negative.

\subsubsection{The $\Psi$-Design Framework}

\begin{center}
\renewcommand{\arraystretch}{1.3}
\begin{tabular}{lll}
\hline
& \textbf{Classical MD} & \textbf{$\Psi$-Design} \\
\hline
Given & Preferences, constraints & Initial $\Psi_0$ \\
Choose & Mechanism/rules & $\Psi$-intervention path \\
Optimize & Outcome at $t$ & $Q$ over time horizon \\
Constraint & Incentive compatibility & $\Psi$-dynamics: $\dot{\Psi} = g(\Psi, \text{mech})$ \\
\hline
\end{tabular}
\end{center}

\paragraph{The $\Psi$-Design Problem.}

\begin{equation}
\max_{\text{mechanism}} \int_0^T Q(\Psi_t) dt \quad \text{s.t.} \quad \dot{\Psi}_t = g(\Psi_t, \text{mechanism})
\label{eq:psi-design}
\end{equation}

This is an optimal control problem where the mechanism affects the state variable $\Psi$.

\subsubsection{The Prediction}

\begin{tcolorbox}[colback=blue!5!white, colframe=blue!75!black, title=Prediction 8: Mechanism-Context Dynamics]
Mechanism design that ignores $\Psi$-dynamics will fail in the long run. 
Mechanisms that achieve short-run efficiency by substituting formal rules 
for social trust will erode the social capital on which future cooperation depends.

\textbf{Formal statement:}
Let $M^*$ maximize static efficiency and $M^{**}$ maximize dynamic welfare.
\begin{equation}
M^* \neq M^{**} \quad \text{when } \frac{\partial \Psi_S}{\partial \Psi_I} < 0
\end{equation}

\textbf{Specific predictions:}
\begin{enumerate}
  \item High-formalization mechanisms show initial success, then gradual erosion
  \item Communities with high initial $\Psi_S$ lose more from formalization
  \item Sustainable mechanisms are ``trust-preserving'': they complement rather than substitute for $\Psi_S$
\end{enumerate}

\textbf{Design principle:} Optimal mechanism satisfies:
\begin{equation}
\frac{\partial \Psi_S}{\partial \text{mechanism}} \geq 0
\end{equation}
i.e., the mechanism preserves or enhances social capital.
\end{tcolorbox}

\subsubsection{Case Studies}

\paragraph{Kidney Exchange: Success.}
\begin{itemize}
  \item Mechanism: Matching algorithm for kidney swaps
  \item $\Psi$-effect: Increases trust in medical system, enables altruistic chains
  \item Why sustainable: Builds on existing altruistic motivation, doesn't substitute for it
  \item Result: Expanding participation, thousands of lives saved
\end{itemize}

\paragraph{High-Stakes Testing: Erosion.}
\begin{itemize}
  \item Mechanism: Standardized tests for school accountability
  \item Initial effect: Performance improves (teaching to test)
  \item $\Psi$-effect: Trust in intrinsic learning motivation erodes (``extrinsic crowding out'')
  \item Long run: Gaming, cheating, teaching to test, loss of love of learning
  \item Result: Initial gains stagnate or reverse
\end{itemize}

\paragraph{Gig Economy Platforms: Mixed.}
\begin{itemize}
  \item Mechanism: Algorithmic matching, ratings, surge pricing
  \item Initial effect: Efficient allocation, flexible work
  \item $\Psi$-effect: Worker-platform trust erodes (algorithmic management)
  \item Long run: Labor organizing, regulatory backlash, platform costs rise
  \item Result: Efficiency gains partially offset by trust costs
\end{itemize}

\subsubsection{Empirical Test Design}

\paragraph{Data.}
\begin{itemize}
  \item Mechanism introductions: Date, type, context
  \item Performance trajectories: Efficiency measures over time
  \item $\Psi$ measures: Trust surveys, informal cooperation indicators
  \item Cross-community comparison: Similar communities, different mechanisms
\end{itemize}

\paragraph{Model.}
Dynamic panel:
\begin{equation}
\text{Performance}_{it} = \alpha_i + \beta_1 \cdot \text{Mechanism}_{it} + \beta_2 \cdot (t \times \text{Mechanism}_{it}) + \gamma \cdot \Psi_{S,it} + \varepsilon_{it}
\end{equation}

The framework predicts $\beta_1 > 0$ (initial positive effect), $\beta_2 < 0$ (erosion over time), 
and the erosion is mediated by $\Psi_S$ decline.

\paragraph{Falsification.}
The prediction is falsified if:
\begin{enumerate}
  \item High-formalization mechanisms show sustained performance without erosion
  \item Social trust ($\Psi_S$) is unaffected by mechanism design
  \item No systematic relationship between formalization and long-run outcomes
\end{enumerate}

\subsubsection{Policy Implications}

\begin{enumerate}
  \item \textbf{Design horizon:} Evaluate mechanisms over long time horizons, not just implementation
  \item \textbf{Trust accounting:} Include $\Psi_S$ effects in cost-benefit analysis
  \item \textbf{Complementary design:} Build mechanisms that enhance rather than substitute for trust
  \item \textbf{Adaptive mechanisms:} Allow mechanism evolution as $\Psi$ changes
  \item \textbf{Exit options:} Preserve informal alternatives to prevent lock-in
\end{enumerate}

%%%%%%%%%%%%%%%%%%%%%%%%%%%%%%%%%%%%%%%%%%%%%%%%%%%%%%%%%%%%%%%%%%%%%%%%%%%%%%%
\subsection{Summary: The Prediction Matrix}
%%%%%%%%%%%%%%%%%%%%%%%%%%%%%%%%%%%%%%%%%%%%%%%%%%%%%%%%%%%%%%%%%%%%%%%%%%%%%%%

\begin{table}[htbp]
\centering
\caption{Summary of Eight Novel Predictions}
\label{tab:predictions-summary}
\small
\renewcommand{\arraystretch}{1.3}
\begin{tabular}{@{}clll@{}}
\toprule
\textbf{\#} & \textbf{Prediction} & \textbf{Key Mechanism} & \textbf{Policy Implication} \\
\midrule
1 & Trade war hysteresis & $\Psi$-dynamics, $\gamma_{S,M} < 0$ & Effects persist post-tariff \\
2 & Cross-domain spillovers & $C^{FS} \neq 0$ & Financial $\to$ political crisis \\
3 & Asymmetric recovery & Convex $\tau(K)$ & Big crises need big response \\
4 & Techlash & $\gamma_{E,S}$, $\gamma_{S,M}$ trade-offs & Technology policy needs $U^S$ \\
5 & Identity non-fungibility & $U^{IDN}$ separate dimension & Money can't buy identity \\
6 & Coherence trap & High-$K$, low-$Q$ stable & Escape needs big push \\
7 & Multi-level conflict & $K^{national} \neq K^{global}$ & Climate needs global $K$ \\
8 & Mechanism-context & $\dot{\Psi} = g(\Psi, \text{mech})$ & Design for $\Psi$-preservation \\
\bottomrule
\end{tabular}
\end{table}

These predictions are:
\begin{itemize}
  \item \textbf{Novel:} Not derivable from any single constituent theory
  \item \textbf{Quantitative:} Grounded in the estimated $\Gamma$-matrix
  \item \textbf{Testable:} With specified falsification criteria
  \item \textbf{Policy-relevant:} With concrete implications
\end{itemize}

They emerge from the unique combinatorial structure of the $(C, \Psi, K, Q)$ 
framework---demonstrating that theoretical integration produces empirical 
payoff beyond what partial theories can achieve.

\vspace{1em}
\hrule
\vspace{0.5em}
\noindent\textit{Cross-references: Chapter~10 ($\Gamma$-matrix estimation), 
Appendix~AN (LLM-MC methodology), 
Appendix~V ($\Psi$-dimensions validation),
Chapter~13 (Policy implications)}

%%%%%%%%%%%%%%%%%%%%%%%%%%%%%%%%%%%%%%%%%%%%%%%%%%%%%%%%%%%%%%%%%%%%%%%%%%%%%%%
